\section{Abordagem Proposta}
\label{sec:metodo}

\subsection{Visão geral do pipeline}
\label{sec:metodo-visao}

\input{figures/architecture}

Os três componentes da Figura~\ref{fig:arquitetura} partem dos turnos de fala da
Seção~\ref{sec:dados-particao}, e os dois ramos que alimentam a simulação resolvem a identidade das
pessoas de formas diferentes. A UDV (Seção~\ref{sec:metodo-udv}) associa o nome citado na matéria aos
cabeçalhos de uma única audiência, e o ramo dos perfis (Seções~\ref{sec:metodo-atores}
e~\ref{sec:metodo-perfis}) associa entre si cabeçalhos de audiências diferentes. Como ambos usam as
mesmas chaves (audiência, índice do turno e offsets na transcrição), o turno que contém a evidência de
uma UDV identifica o registro de ator correspondente, e dizemos que a opinião está ligada a esse ator.
É por esse cruzamento que uma opinião publicada sobre uma audiência de validação ou de teste se torna
uma pergunta sobre um ator com perfil (Seção~\ref{sec:setup}).

\subsection{Unidades Deliberativas Verificáveis}
\label{sec:metodo-udv}

As opiniões do PublicHearingBR registram o que a matéria atribui a cada pessoa, mas o conjunto não
indica que trecho da transcrição corresponde a cada opinião (Seção~\ref{sec:dados}). Sem essa
indicação, não é possível apresentar ao leitor a passagem da fala a que a opinião se refere nem
verificar se a fala a sustenta. A UDV estabelece o vínculo entre opinião e fala com um registro por
opinião que aponta, quando existe, o trecho da transcrição usado como evidência
(Figura~\ref{fig:udv}).

\textbf{Resolução da pessoa.} O resumo estruturado registra o nome usado na matéria, que costuma
diferir do cabeçalho da transcrição (``Rodrigo Marinho'' na matéria, ``RODRIGO SARAIVA MARINHO'' na
transcrição), e quem preside a sessão aparece como \textit{PRESIDENTE}, com o nome entre parênteses.
Cada turno fornece como candidatos o nome do cabeçalho e, quando há, o nome entre parênteses, que é o
de quem preside nos turnos de presidência e pode ser um nome social nos demais. Um participante é
associado a um candidato quando os dois conjuntos de palavras, sem acento e em maiúsculas, são iguais
ou um contém o outro, desde que ambos tenham pelo menos duas palavras. Um nome formado por uma
palavra é associado a um candidato idêntico ou, na falta deste, ao único candidato da audiência que
contenha essa palavra; havendo mais de um, nenhum turno é associado à pessoa. O participante recebe
os turnos de todos os candidatos a que foi associado. Com essa regra, 1.020 dos 1.065 participantes
são associados a pelo menos um turno, e os 45 restantes respondem por 90 opiniões; a precisão da
associação não foi medida.

\input{figures/udv-example}

\textbf{Sentenças candidatas.} A segmentação em sentenças é feita dentro de cada turno da pessoa, e
cada sentença guarda o índice do turno de origem. Com isso, nenhuma sentença junta o fim de um turno
ao início do turno seguinte da mesma pessoa, entre os quais a transcrição registra falas de outros
participantes. São mantidas as sentenças com quatro palavras ou mais que não se reduzam a
marcação de palco, como ``(Palmas.)'', o que resulta em 115.599 sentenças candidatas nas 206
audiências.

\textbf{Citação literal.} A matéria apresenta como fala literal os trechos entre aspas das opiniões,
o que permite um vínculo por casamento de texto, sem modelo, mas pode editá-los. Na
Figura~\ref{fig:udv}, ela escreve ``presidente da comissão, que é eleita'', e a transcrição registra
``presidente que foi eleita''. Por isso, a busca usa prefixos da citação e não exige que o restante
dela coincida com a fala.
Das 2.113 opiniões de pessoas associadas a turnos, 960 têm alguma citação com 10 caracteres ou mais.
De cada citação são extraídos prefixos de 10, 6, 4 e 3 palavras (a citação inteira, quando é mais
curta que o prefixo), buscados nessa ordem em cada turno da pessoa, sem distinção entre maiúsculas e
minúsculas e respeitando as fronteiras de palavra. Entre as citações de uma mesma opinião, prevalece a
de prefixo encontrado mais longo e, no empate, a primeira. Quando esse prefixo tem seis palavras ou
mais, a citação é considerada confiável, e a sentença que o contém é a evidência, com o nível
\texttt{quote\_found}; se o prefixo atravessa o fim de uma sentença, a evidência reúne as sentenças
que ele toca. Se o prefixo ocorre mais de uma vez na fala, é escolhida a sentença com maior índice de
Jaccard de palavras em relação à opinião.

Prefixos de três ou quatro palavras costumam ser expressões comuns, encontradas em pontos da fala sem
relação com a citação. O limite de seis palavras foi definido a partir de uma leitura exploratória
dos casamentos agrupados por tamanho de prefixo, sem protocolo de anotação. Das 960 opiniões com
citação, 549 têm algum prefixo encontrado na fala da pessoa e 277 têm prefixo de seis palavras ou
mais.

\textbf{Similaridade semântica.} Quando a opinião não tem citação confiável, ela e as sentenças
candidatas da pessoa são codificadas pelo Serafim 335m~\citep{gomes2024serafim}, um codificador de
sentenças para o português (vetores de 1.024 dimensões, entrada limitada a 128 tokens). A evidência é
a sentença de maior similaridade de cosseno com a opinião. Um codificador denso foi preferido ao
TF-IDF porque representa o sentido da sentença em contexto, e a matéria costuma parafrasear a fala,
caso em que a comparação por vocabulário compartilhado falha. Se o prefixo mais longo encontrado tem
menos de seis palavras (prefixo curto) e a sentença que contém a sua primeira ocorrência na fala coincide com a
escolhida pelo codificador, a UDV registra a coincidência como corroboração fraca, sem alterar o
nível. Cada UDV registra a versão do codificador e o corte com que foi produzida.

\textbf{Nível e proveniência.} O nível resume o tipo de vínculo encontrado
(Tabela~\ref{tab:udv-niveis}), e os níveis \texttt{semantic\_match\_high} e
\texttt{semantic\_match\_weak} indicam de que lado do corte ficou o cosseno da melhor sentença, sem
afirmar que ela sustenta a opinião. A proveniência registra a origem do vínculo: \texttt{weak}
(rótulo produzido por regra, sem conferência humana) para o casamento de texto e \texttt{model} para
o codificador. Nenhum registro tem proveniência humana.

\textbf{Localização e cobertura.} A sentença escolhida é localizada na transcrição original por uma
expressão regular restrita ao turno de origem, que devolve a primeira ocorrência do texto no turno, e
todas as 2.105 evidências foram localizadas com offsets de caracteres. As 98 opiniões restantes não
têm evidência. Nos 8 registros \texttt{no\_evidence}, a única fala transcrita da pessoa é a marcação
``(Manifestação em LIBRAS.)'', e entre as pessoas sem turno associado há casos de grafia divergente do
nome, de instituições citadas como participantes e de pessoas que não falaram na sessão. Nenhum
desses níveis indica erro da matéria.

\input{tables/udv-tiers}

\textbf{Calibração do corte.} O corte que separa \texttt{semantic\_match\_high} de
\texttt{semantic\_match\_weak} foi calibrado nas audiências de treino, porque a amostra de validação
humana é sorteada no teste e depende dessa fronteira. A regra declarada como principal escolheria o
corte que maximizasse o índice J de Youden~\citep{youden1950index} na separação entre acertos e erros
do codificador. As consultas dessa regra são opiniões com citação confiável das quais os trechos entre
aspas foram removidos, e há acerto quando a sentença mais similar à opinião é a da própria citação.
Das 183 opiniões do treino com citação confiável, 142 mantêm quatro palavras ou mais após a remoção;
nas 41 restantes, a opinião se reduz quase inteiramente à citação. A regra se mostrou degenerada: só
12 das 142 consultas são acertos, os erros têm cosseno mais alto que os acertos (médias de 0,655 e
0,578, respectivamente), e o maior valor de J é zero, obtido por um corte que aceita todas as
consultas.

O corte adotado depois dessa falha vem de uma regra de pares que a declaração da regra principal
previa só para comparação, com as 183 opiniões do treino com citação confiável como positivos, cada uma pareada com a
sentença da própria citação. Os negativos são a mesma opinião pareada com uma sentença sorteada de
outra audiência do treino. O corte é o ponto médio entre o primeiro quartil dos positivos (0,628) e
o terceiro quartil dos negativos (0,281), igual a 0,454 e arredondado para 0,45 (IC 95\% por
\textit{bootstrap} de audiências, de 0,437 a 0,469). Com negativos sorteados da fala da própria
pessoa, que correspondem à escolha que o codificador de fato realiza, a mesma regra resultaria em
0,50, e a escolha entre as duas variantes não foi validada. Como 43 das 1.828 UDVs semânticas (2,4\%)
estão abaixo de 0,45, o nível \texttt{semantic\_match\_weak} reúne poucos registros, e o corte não
foi validado como medida de confiança.

\textbf{Verificação.} Um verificador separado refaz, a partir das transcrições e dos resumos
estruturados, a resolução de pessoa, a extração das sentenças candidatas e a busca de citações de
cada opinião. Em seguida, confere se a opinião e a pessoa coincidem com o resumo estruturado, se a
evidência está num único turno da pessoa e se os offsets recuperam o texto registrado. Também confere
se o nível é coerente com o resultado da busca de citações, o cosseno e o corte. Como reutiliza os procedimentos da
construção nessas três etapas, o verificador confere a consistência dos registros com esses
procedimentos, e não a correção deles. Ele
também não recalcula a escolha do codificador, não confere se a sentença se repete no turno e não
avalia se a evidência sustenta a opinião; no conjunto descrito aqui, não acusou nenhuma
inconsistência.

\subsection{Falas por ator entre audiências}
\label{sec:metodo-atores}

Como a UDV trata uma audiência de cada vez, descrever o que uma pessoa defende ao longo do tempo
exige reunir suas falas em todas as audiências de que participou e, para isso, decidir quando
cabeçalhos de audiências diferentes se referem à mesma pessoa. A regra de nome contido da UDV compara
nomes dentro de uma única audiência, que tem poucos falantes; aplicada entre audiências, ela uniria
pessoas diferentes, como ``Luiz Lima'' e ``Gustavo Luiz de Lima Correia''. Por isso, o registro de ator usa como chave o
nome do cabeçalho sem acento e em maiúsculas (nos turnos de presidência, o nome entre parênteses),
complementado por uma lista de mesclas.

A lista de mesclas foi proposta a partir dos 153 pares de nomes candidatos, em que um nome contém o
outro ou os dois têm as mesmas palavras, e validada com evidências do próprio conjunto. Essas
evidências são o cargo informado na matéria, a apresentação feita pelo presidente da sessão, a
autoapresentação, o partido e a UF do cabeçalho e a data da audiência. O resultado tem 37 grupos de mescla, uma chave
dividida por audiência (um caso de homônimo exato) e um par sem resolução, mantido separado. Nenhuma
identidade foi conferida em registros externos ao conjunto.

Três regras de filtragem retiram os turnos que não trazem a fala da própria pessoa ou que se limitam
à condução da sessão. São removidos os turnos de intérprete e de tradução, cuja fala pertence a outro
participante, os turnos de presidência com menos de 50 palavras e os turnos vazios ou formados por
marcação de palco. Nas 206 audiências, os turnos de presidência concentram 42,7\% das palavras dos
falantes presentes em duas ou mais audiências, contadas antes dos filtros e das mesclas, e boa parte
dessa fala é de condução da sessão. Descartá-los por completo,
porém, eliminaria opiniões, pois 336 das 2.105 evidências de UDV estão em turnos de presidência. O
turno de presidência mediano tem 28 palavras, e o limite de 50 palavras mantém 2.264 dos 6.266 turnos
de presidência, com 89\% de suas palavras, e 332 das 336 evidências.

Após os três filtros restam 13.190 dos 17.264 turnos, com o texto copiado da transcrição sem limpeza
e conferido com o intervalo indicado pelos offsets. Cada registro de ator traz o nome, o partido e
a UF do cabeçalho e, por audiência, os turnos com índice, papel (presidência ou fala comum), offsets e
texto. Os 301 atores que falam em duas ou mais audiências formam o material dos perfis, porque neles
há fala de audiências diferentes, o que permite escrever o perfil com parte das audiências e avaliá-lo nas
demais. Desses, 264 falam em pelo menos uma audiência de treino e recebem perfil
(Seção~\ref{sec:dados-particao}), 79 deles com uma única audiência de treino.

\subsection{Perfis por ator}
\label{sec:metodo-perfis}

A fala de um ator se distribui por audiências de temas diferentes e se mistura a passagens de
condução da sessão e de cortesia. O perfil condensa essa fala em um texto escrito por um LLM e
destinado a servir de instrução a um modelo que deve responder como a pessoa responderia, tanto sobre
temas que ela já tratou quanto sobre temas novos. Por isso, além do que a pessoa defende, critica,
propõe e cobra, o perfil registra os critérios que ela invoca para justificar essas posições, o que
ela declara sobre si e sobre o próprio lado e a forma como argumenta. Na geração final, cada perfil é
escrito com as audiências de treino (Seção~\ref{sec:dados-particao}), para que nenhuma fala ou
opinião das audiências de validação e de teste chegue ao prompt.

Cada ator recebe uma única geração, cujo prompt do usuário contém todas as suas falas de treino,
organizadas em blocos por audiência e em ordem cronológica. Cada bloco é aberto por um cabeçalho com a
data da matéria e o assunto registrado no resumo estruturado, e os turnos de presidência são
precedidos da marca \texttt{[presidência da sessão]}. A data permite situar no tempo eventuais
mudanças de posição, e o assunto ajuda a resolver referências vagas, como ``este projeto de lei''. O
prompt de sistema impõe quatro grupos de regras, e os dois prompts estão no
Apêndice~\ref{ap:prompt}.

As regras de fonte única restringem o perfil às falas, sem conhecimento prévio sobre a pessoa, os
partidos ou os temas em debate. Nada que apareça só no cabeçalho é atribuído à pessoa, propostas de
terceiros que ela relata ou questiona não se tornam posições dela, e siglas, nomes e datas não são
completados. As de condução e cortesia determinam que sejam desconsideradas passagens como
passar a palavra, controlar o tempo e agradecer. As de blocos organizam o perfil em até quatro blocos (posições;
critérios e valores; alinhamentos declarados; forma de argumentar), omitem o bloco sem apoio nas falas
e proíbem deduzir valores ou alinhamentos das posições ou do partido. Na forma de argumentar, um traço
precisa aparecer em mais de um turno sem ser contrariado. As de escrita pedem itens ordenados por
importância, proíbem descrever o que a pessoa não disse e fixam um teto de 800 palavras, que não é
imposto durante a geração.

A divisão em blocos destina-se a separar o material que orienta respostas sobre temas já tratados
(as posições) do material que orienta respostas sobre temas novos (os critérios, os alinhamentos
declarados e a forma de argumentar). As proibições de inferência visam impedir que o perfil atribua à pessoa
características do partido ou das posições que ela própria não enunciou: com um alinhamento deduzido
do partido, o modelo que recebe o perfil responderia com base nele, e não nas falas.
Pelo mesmo motivo, o perfil não descreve ausências, porque uma frase como ``não menciona temas de
saúde'' poderia ser tomada como posição.

As regras sobre o cabeçalho, sobre posições de terceiros, sobre siglas e nomes e sobre o bloco de
forma de argumentar foram acrescentadas após a leitura de perfis gerados durante o desenvolvimento do
prompt. Esses perfis foram gerados por outros LLMs e, em parte, com falas de todas as audiências,
inclusive as de teste. Neles, o assunto da audiência aparecia como posição do ator, a proposta de um
convidado resumida pelo presidente aparecia como proposta do próprio presidente, e siglas e sobrenomes
eram completados com conhecimento externo às falas. Além disso, uma frase dirigida a uma única
autoridade aparecia como traço geral da pessoa, e um traço atribuído a ela era contrariado por falas
de outra audiência. As regras resultantes tratam da forma do perfil, e não do conteúdo de posições
específicas, e os perfis de desenvolvimento não entram na avaliação.

Os dois prompts compõem uma conversa de duas mensagens, formatada pelo \textit{chat template} do
próprio LLM, sem formatação adicional. A geração usa amostragem com temperatura 0,2, \textit{top-p}
de 0,9 e limite de 3.000 tokens de saída, e a semente é fixada antes de cada geração, para que o
perfil de um ator não dependa da ordem em que os atores são processados. Uma resposta que atinge o
limite de tokens sem terminar é considerada falha e descartada. O modelo gerador é <modelo>, e o
prompt é enviado inteiro, sem truncamento; o maior tem <n> tokens de entrada, <dentro ou fora> da
janela de contexto desse modelo. Cada perfil registra o modelo, a versão dos prompts, as audiências
usadas e o número de tokens de entrada e de saída.

% Pendente: trecho de um perfil real (Quadro) depois da rodada de geração com o modelo final.

\subsection{Simulação de atores}
\label{sec:metodo-simulacao}

Na simulação, o perfil é fornecido a um LLM, o modelo simulador, que deve responder como a pessoa
responderia, escolhendo entre opiniões (múltipla escolha) ou escrevendo uma fala (geração aberta); nos
experimentos, esse modelo é o mesmo LLM que gera os perfis. Três elementos (o perfil, os trechos
recuperados e a \textit{classifier-free guidance}, CFG) são acrescentados um de cada vez, o que define
quatro condições encadeadas. A condição 0 recebe só o nome e o cargo da pessoa; a 1 acrescenta o
perfil; a 2, trechos das falas de treino recuperados pelo assunto; e a 3 combina a 2 com a 0 por
CFG.
% Pendente: motivo de usar o mesmo LLM para gerar os perfis e para simular.

\textbf{Prompt comum.} Todas as chamadas partem do mesmo prompt de sistema e diferem só no material
incluído, porque, com instruções diferentes por condição, a diferença entre condições vizinhas e o
contraste da CFG mediriam também a troca de instrução (Apêndice~\ref{ap:prompt-simulacao}). O prompt
apresenta a pessoa pelo nome e pelo cargo: na múltipla escolha, o cargo registrado na matéria da
audiência avaliada; na geração aberta, o da audiência de treino mais recente em que a pessoa tem
opinião ligada, ou nenhum. As instruções proíbem o modelo de usar conhecimento prévio sobre a
pessoa, os partidos e o grupo que ela representa, assim como a própria opinião. Quando o material não
trata do assunto, as instruções orientam o modelo a se basear nos critérios e nos alinhamentos da
pessoa. O partido não é fornecido em nenhuma condição, para que a falta de material não seja suprida
pela posição associada a ele. A sigla entre parênteses, com ou sem a UF, é retirada do cargo
(``Deputado (PT-SP)'' torna-se ``Deputado''), e o partido só chega ao modelo quando a própria pessoa o
declara nas falas. Cada item do perfil no formato pedido recebe um identificador formado pela inicial do bloco e pela ordem do item (P1, C1, A1, F1), e na condição 0 o
perfil é substituído por ``Nenhum.''.

\textbf{Trechos recuperados.} Como o perfil condensa as falas de treino num texto com teto de 800
palavras, detalhes que podem decidir uma pergunta são omitidos. Para recuperá-los, o assunto da
audiência é codificado pelo Serafim 335m da UDV e comparado às sentenças de quatro palavras ou mais
dos turnos da pessoa nas audiências de treino. Entram as $k$ sentenças de maior cosseno, no máximo
uma por turno, cada uma com a sentença anterior e a seguinte do mesmo turno. Os trechos são incluídos
na mensagem do usuário em ordem de data, para que uma mudança de posição apareça na ordem em que
ocorreu. Cada um traz um identificador (T1 a T$k$), a data e o assunto da audiência de origem e,
quando for o caso, a marca de presidência; o cosseno não é mostrado ao modelo.

\textbf{\textit{Classifier-free guidance}.} Mesmo com o perfil no prompt, o modelo pode responder com a
própria opinião sobre o tema ou com o que associa ao nome da pessoa. A
CFG~\citep{ho2022classifier,sanchez2024stay} amplifica, a cada passo da geração, a diferença entre um
prompt completo $x$ e um prompt negativo $\bar{x}$:
\begin{equation}
  \tilde{p}(y_t \mid y_{<t}) \propto p(y_t \mid \bar{x}, y_{<t})
  \left[ \frac{p(y_t \mid x, y_{<t})}{p(y_t \mid \bar{x}, y_{<t})} \right]^{\gamma},
  \label{eq:cfg}
\end{equation}
em que $\gamma = 1$ reproduz a geração com $x$ e $\gamma > 1$ desloca a distribuição na direção
indicada pelo material presente em $x$ e ausente em $\bar{x}$. O prompt negativo é o da condição 0,
com o perfil ``Nenhum.'', sem trechos e sem exemplos. Na múltipla escolha, a condição 3 aplica a
Equação~\ref{eq:cfg} às probabilidades das letras já calculadas nas condições 0 e 2, sem passagem
adicional pelo modelo. Restrita às quatro letras, a combinação pondera cada opção por
$p_2^{\gamma}\, p_0^{1-\gamma}$ e, com $\gamma > 1$, penaliza as opções que a condição 0 já favorece.
Ela é aplicada em cada rotação, antes da média entre rotações (Seção~\ref{sec:setup}), porque
combinar depois da média produz outro resultado. Na geração aberta, a combinação é feita a cada token da fala, com duas passagens pelo
modelo por token.

\textbf{Pedidos.} Além da múltipla escolha, usada na avaliação (Seção~\ref{sec:setup}), a simulação
tem três pedidos, todos com decodificação gulosa: a base da estimativa, a fala e a justificação
(Apêndice~\ref{ap:prompt-simulacao}). Na base da estimativa, o modelo classifica o material como
DIRETA, INDIRETA, ESPECULATIVA ou SEM BASE, em ordem decrescente de proximidade com o assunto. A classificação
é feita numa chamada própria e sem CFG porque o prompt negativo, sem material, pode favorecer a
recusa, e, com $\gamma > 1$, a CFG afastaria a resposta da recusa justamente quando ela é devida. Com
SEM BASE, a fala não é gerada, e a saída é ``o material não permite estimar''; nos demais casos, o
modelo escreve em primeira pessoa uma fala sobre o assunto e, numa chamada separada e sem CFG, a
justificação dessa fala.

Na geração aberta, o prompt de sistema traz também até dois exemplos fixos de fala da pessoa, de 60 a
150 palavras, sorteados com semente fixa entre os turnos de treino que não são de presidência e,
quando possível, de audiências diferentes. Cada exemplo começa na segunda sentença do turno, para
excluir os cumprimentos de abertura. Os exemplos se destinam a mostrar o estilo da pessoa e são
excluídos da múltipla escolha, porque, como contêm posições, fariam a diferença entre as condições 0
e 1 medir também o efeito dessas falas.

\textbf{Justificação conferida.} Na justificação, o modelo recebe a fala simulada e escreve, para cada
frase, um objeto JSON com os identificadores do material que a apoia e, para cada exemplo ou trecho
citado, a parte copiada. A conferência é automática: cada identificador precisa existir no material
da chamada, e cada cópia é aceita quando as suas seis primeiras palavras aparecem no exemplo ou trecho
citado, pela regra de prefixo da citação literal (Seção~\ref{sec:metodo-udv}). Uma frase é
considerada sem fundamento quando não aparece em nenhum objeto ou quando algum objeto que a contém não
tem apoio, cita identificador inexistente ou traz cópia ausente, com menos de seis palavras ou não
encontrada. A conferência mostra que os itens citados existem e que o início de cada cópia é literal,
mas não que o item sustenta a frase nem que o modelo o usou, pois a justificação é produzida depois
da fala.
